% Part: computability
% Chapter: computability-theory
% Section: computable-sets

\documentclass[../../../include/open-logic-section]{subfiles}

\begin{document}

\olfileid{cmp}{thy}{cps}
\olsection{સંગણનીય ગણો}

સંગણનીયતાની કલ્પનાને સંગણનીય વિધેયોથી સંગણનીય ગણો સુધી
વિસ્તારી શકાય છે:

\begin{defn}
  માનો કે $S$ પ્રાકૃતિક સંખ્યાઓનો ગણ છે. તો $S$ ત્યારે અને માત્ર
  ત્યારે \emph{સંગણનીય} છે જ્યારે તેનું લક્ષણ વિધેય~$\Char{S}$
  સંગણનીય હોય. બીજા શબ્દોમાં, $S$ ત્યારે અને માત્ર ત્યારે
  સંગણનીય છે જ્યારે
\[
\Char{S}(x) =
\begin{cases}
1 & \text{જો $x \in S$} \\
0 & \text{અન્યથા}
\end{cases}
\]
વિધેય સંગણનીય હોય. એ જ રીતે, સંબંધ $R(x_0, \dots, x_{k-1})$
ત્યારે અને માત્ર ત્યારે સંગણનીય છે જ્યારે તેનું લક્ષણ વિધેય
સંગણનીય હોય.

સંગણનીય ગણો અને સંબંધોને \emph{નિર્ણેય} પણ કહે છે.
\end{defn}

\begin{explain}
હવે સંગણનીયતાની ઘણી કલ્પનાઓ છે: આંશિક વિધેયો માટે, વિધેયો
માટે અને ગણો માટે. તેમને ગૂંચવશો નહીં!{} \iftag{TMs}{આંશિક વિધેય
  ગણતું ટ્યુરિંગ મશીન, વિધેય જ્યાં વ્યાખ્યાયિત છે એવા ઇનપુટ માટે,
  વિધેયનું આઉટપુટ આપે છે; ગણ નક્કી કરતું ટ્યુરિંગ મશીન ઇનપુટ
  તે ગણમાં છે કે નહીં તે નક્કી કરીને $1$ અથવા~$0$ આપે છે.}{}
\end{explain}

\end{document}
